\documentclass[journal]{IEEEtran}
%------------------------------------------------------------------------------
% IEEE Transactions on Geoscience and Remote Sensing (TGRS) submission
% Converted from the Springer Nature sn-jnl manuscript (sn-article.tex).
% Non-essential content lives in the online supplement (tgrs-supplement.tex):
% it is referenced here as "Supplementary Material, Sec./Table S#/Fig. S#".
%------------------------------------------------------------------------------
\usepackage{amsmath,amssymb,amsfonts}
\usepackage{graphicx}
\graphicspath{{figures/}}
\usepackage{booktabs}
\usepackage{array}
\usepackage{xcolor}
\usepackage{url}
\usepackage[hidelinks]{hyperref}

\begin{document}

\title{Deep Architectures Can Generalize: Sparse Feature Selection Improves Spatial Transfer in Earth Observation}

\author{Michael~L.~Mann,~Sairam~Venkatachalam,~Disha~Kacha,~Devarsh~Sheth,~Ryan~Engstrom,~and~Amir~Jafari
\thanks{All authors are with the Department of Geography \& Environment and the Data Science Program, The George Washington University, Washington, DC 20052 USA.}
\thanks{Corresponding author: Michael L. Mann (e-mail: mmann1123@email.gwu.edu).}
\thanks{This research received no external funding.}}

\markboth{IEEE Transactions on Geoscience and Remote Sensing}%
{Mann \MakeLowercase{\textit{et al.}}: Sparse Feature Selection Improves Spatial Transfer in Earth Observation}

\maketitle

\begin{abstract}
Accurate crop classification underpins food-security monitoring across the dryland Global South, yet operational mapping there means applying a model to regions where no labels were collected. Most studies instead validate within a single region (random or field-wise cross-validation), which measures in-distribution fit rather than spatial transfer. On a Western Cape, South Africa Sentinel-2 dataset (optical only), we benchmark twenty-three classical, deep, and patch-based models twice: in-region (field-wise) and on a spatially disjoint holdout tile. The rankings invert. In-region the dense temporal deep nets lead (macro-F1 0.72--0.78); under transfer this cluster fractures, as L-TAE and a Transformer encoder drop from 0.78 to 0.58 and CNN-BiLSTM falls the most (0.72 to 0.46), while a TabNet ensemble becomes the best model (0.60) and simple trees and logistic regression (0.56) prove the most robust. The cause is mechanistic: sparse, axis-aligned feature selection is what survives transfer; TabNet inherits it through sparsemax masks, while dense nets discard it. A control isolates it: a Transformer with far greater attention capacity transfers no better than L-TAE, so capacity is not the lever; sparsity is. We graft that bias onto a deep network, L-TAE-S, which ties TabNet for the best transfer with the smallest in-region-to-holdout gap of any model. Field-level aggregation is a second, substitutable lever, and robust models hold near-full accuracy on a quarter of the fields. The message is twofold: model selection in crop mapping is an artifact of the validation protocol, and the property deciding whether a model survives transfer is sparse, axis-aligned feature selection.
\end{abstract}

\begin{IEEEkeywords}
Crop classification, deep learning, domain generalization, feature selection, remote sensing, Sentinel-2, spatial transfer, time-series classification.
\end{IEEEkeywords}

\section{Introduction}
\IEEEPARstart{A}{ccurate} and timely crop classification is foundational to precision agriculture, enabling yield forecasting, resource optimization, and food-security assessment \cite{begue2020, ieee2017}. The Sentinel-2 constellation, with its multispectral richness and roughly five-day revisit, has made dense multi-temporal optical imagery widely available. Yet reliable crop classification remains hard throughout Africa and the Global South: spectral similarity between crops, irregular planting and phenological cycles, persistent cloud cover, and rain-fed dryland management all degrade accuracy \cite{begue2020, mahlayeye2024}. Classical methods such as Random Forest (RF) and XGBoost have proven effective and efficient on engineered features \cite{mahlayeye2024, mann2023, mann2017, mann2019}, but depend on feature engineering and can be inconsistent across regions \cite{maxwell2019, belgiu2016}. Deep learning---recurrent and convolutional temporal models---reduces reliance on hand-crafted predictors and captures phenological structure that summary statistics miss \cite{zhu2017, russwurm2018, tempcnn2019}, but its generalizability in complex, low-label settings remains uncertain \cite{pankajakshan2025}.

A second, and for our purposes critical, gap concerns \emph{how generalization is measured}. Much of the crop-classification literature reports accuracy under random or field-wise cross-validation drawn from a single region \cite{belgiu2016, ltae2020, tempcnn2019, russwurm2018, ieee2017}. Such protocols hold out individual pixels or fields but not \emph{space}: training and test observations share the same atmosphere, soils, management calendars, and acquisition geometry, so the reported numbers reflect in-distribution fit rather than transfer to new terrain. Because operational maps are produced by applying a model to regions where no labels were collected, this distinction matters, yet studies that re-rank models under a spatially disjoint holdout remain rare \cite{pankajakshan2025, jin2019, rustowicz2019}. This differs from the bulk of the transfer-learning literature, which adapts a model using target-domain data---labels, unlabeled imagery, or self-supervised pre-training \cite{ma2024, cropformer2023, mai2025}. Our setting is the no-target-access case (domain generalization): models are trained on source tiles and applied to the holdout directly, with no adaptation, so what determines transfer is a property intrinsic to the model rather than a target-specific tuning step.

The Western Cape of South Africa concentrates the operational challenges---smaller fields, rain-fed dryland systems, heterogeneous planting rotations, and persistent winter cloud cover---while offering an unusually rich, high-quality labeled dataset \cite{mlhub}. This makes it a realistic test bed: rich enough to train and compare modern architectures fairly, yet hard enough to expose the failure modes that matter in practice. We evaluate CNN-BiLSTM hybrids, TabNet ensembles, lightweight temporal attention (L-TAE) and full Transformer-encoder networks, temporal-convolution (TempCNN) networks, and patch-based 2D/3D CNNs, benchmarked against XGBoost, LightGBM, and RF. Critically, we evaluate every model twice: under field-wise cross-validation \emph{within} the training region, and on a spatially disjoint holdout tile that no training pixel touches. We restrict inputs to multi-temporal Sentinel-2 optical imagery alone, to isolate the spectral-time-series signal and maximize applicability where only optical data are available. All code is public \cite{primary_repo}.

This study makes five contributions. First, a multi-architecture benchmark (twenty-three classical, deep, and patch-based models) showing that in-region versus spatial-transfer evaluation produce \emph{inverted} rankings on the same data. Second, a mechanistic account of \emph{why}: sparse, axis-aligned feature selection, the property behind two decades of RF and gradient-boosted-tree dominance, is what survives the move to a holdout tile, and TabNet inherits it via sparsemax masks. Third, a controlled architectural comparison isolating sparse feature selection rather than attention capacity as the lever, including L-TAE-S, a sparse-attention encoder we introduce that ties TabNet for the smallest transfer gap of any model. Fourth, identification of field-level aggregation as a second, substitutable robustness lever that recovers transfer for dense models but is redundant for sparse ones. Fifth, a training-set reduction experiment and the first evaluation of \texttt{xr\_fresh} automated time-series feature extraction in crop classification. This leads to our central question: across machine-learning, deep-learning, and hybrid architectures, which most reliably classifies crops from multi-temporal Sentinel-2 optical imagery \emph{when evaluated on a spatially disjoint region}, and what model properties determine whether in-distribution accuracy survives the transfer?

\subsection{Related Work}
Early crop classification relied on classical learners (RF, SVM, logistic regression, gradient-boosted trees) operating on engineered descriptors---spectral indices, temporal summaries---valued for efficiency, interpretability, and robustness on structured data \cite{maxwell2019, belgiu2016, breiman2001, chen2016, ke2017}, though their performance depends on feature engineering that can limit transferability across agroecological regions \cite{wardlow2007towards, mahlayeye2024}. Deep models learn representations directly from reflectance sequences: CNNs capture local spectral-spatial pattern \cite{cnn_multitemp}, and recurrent and attention architectures model seasonal phenology \cite{ieee2017, russwurm2018, ltae2020, russwurm2020}. Hybrid designs combine convolutional feature extraction with temporal modeling \cite{fan2023hcrnn, tufail2025cnnrnn, bandar2024bilstm}, or use deep backbones as feature extractors feeding classical classifiers \cite{espejo2020, koklu2022, yang2020}. A recurring assumption is that richer learned representations generalize better than compact engineered ones; whether this holds under genuine spatial transfer, rather than within-region validation, is less established and is a central question here. Architectures that explicitly target cross-scenario generalization (e.g., Cropformer \cite{cropformer2023}) remain the exception and are rarely benchmarked against classical and deep baselines on a common spatially disjoint holdout. We fill this gap with a systematic comparison under \emph{both} protocols, and additionally evaluate \texttt{xr\_fresh} as an automated time-series feature-extraction framework in a crop-classification setting \cite{mann2025}.

\section{Data}
The study uses a labeled field dataset distributed through Radiant Earth's MLHub \cite{mlhub} and used as the basis of the AI4FoodSecurity (South Africa) challenge \cite{ai4foodsecurity}, in an agriculturally productive part of the Western Cape. Two adjacent Sentinel-2 tiles (\texttt{34S\_19E\_258N}, \texttt{34S\_19E\_259N}) form the training region (4{,}150 labeled fields over 622~km\textsuperscript{2}; mean field 15.0~ha); a third, immediately adjacent tile (\texttt{34S\_20E\_259N}) is reserved as a holdout for spatial-transfer evaluation (2{,}417 fields over 282~km\textsuperscript{2}). The tiles do not overlap, so no pixel or field is shared, and transfer must contend with different soils, geometries, and acquisition conditions. Because the holdout lies within the same agroecological zone and growing calendar, this is a relatively \emph{local} spatial shift, and the gaps we report are best read as a conservative lower bound on degradation under broader transfer. Throughout, ``in-region'' denotes a field-wise train/validation/test split within the training region; ``spatial-transfer'' (ST) or ``holdout'' denotes predictions on the disjoint tile.

We classified five crop types: wheat, canola, lucerne/medics (alfalfa), barley, and small-grain grazing. The classes are heavily imbalanced---lucerne/medics is 43\% of training fields and 56\% of the holdout (Supplementary Material, Fig.~S1)---which motivates macro-F1 as the primary metric (Sec.~\ref{sec:metrics}) and class-imbalance handling during training. Spectrally, wheat and barley overlap strongly, the dominant source of persistent confusion (Supplementary Material, Fig.~S4). Sentinel-2 Level-2A surface reflectance was acquired for the January--December 2017 growing season, cloud-masked with \texttt{s2cloudless} plus a shadow mask, and composited monthly. Six features per month are used---bands B2, B6, B11, B12, plus EVI and Hue, with EVI and Hue computed per scene before compositing---chosen to span the most informative spectral regions while limiting redundancy among correlated bands \cite{mann2023, tufail2025cnnrnn}. June 2017 yielded no usable imagery and May was discarded for cloud, leaving 10 monthly observations, so each pixel is a 60-dimensional time series (6 features $\times$ 10 months). Field polygons were digitized from sub-meter aerial photography and carry a single seasonal crop label \cite{mlhub}; all predictions are made at the field level.

\section{Methodology}
We compare three paradigms: feature-based classical machine learning, representation-learning deep networks, and patch-based spatial models. To prevent leakage, all train/validation/test splitting is performed at the field level with a fixed random seed, so no field's pixels straddle partitions; pixel- and patch-level predictions are aggregated to a field label by majority vote (or, for probabilistic models, mean pooling). Computational cost across paradigms is reported in the Supplementary Material (Sec.~S-C).

\subsection{Classical Machine Learning}
We automate feature extraction with the \texttt{xr\_fresh} toolkit, computing a set of statistical and temporal descriptors (e.g., autocorrelation, linear trend, day-of-year of the seasonal maximum, quantiles, complexity) from each pixel's 10-month profile, yielding a 174-dimensional per-pixel vector \cite{mann2023, mann2025}; the full feature list is in the Supplementary Material (Sec.~S-A, Table~S1). Four baselines---logistic regression, RF (300 trees), LightGBM, and XGBoost \cite{james2013, breiman2001, ke2017, chen2016, friedman2001}---are trained on this representation, all of them tree- or margin-based with built-in feature selection or shrinkage; standard formulations are reproduced in the Supplementary Material (Sec.~S-A). We also report soft-voting and stacked ensembles (RF/XGBoost/LightGBM, with logistic-regression meta-learner), an Optuna-tuned field-level XGBoost, and a SMOTETomek-resampled variant; the \emph{Stacking} and \emph{SMOTE Stacked} pair differs only in resampling, isolating the effect of class rebalancing. At the field level, pixel features are mean-pooled within each polygon before classification \cite{penabarragan2011}.

\subsection{Deep Learning Architectures}
At the pixel level each observation is a 60-feature temporal sequence. Although the pixel models train on millions of pixels, pixels within a field are highly autocorrelated and share one label, so the number of statistically independent units is closer to the number of fields; field-wise splitting keeps evaluation honest at the field scale. We evaluate, as five-seed ensembles, a CNN-BiLSTM hybrid trained with weighted focal loss; TabNet \cite{tabnet_orig, tabnet2024}, which treats the observations as tabular input and, at each decision step, uses an attentive transformer to produce a \emph{sparse} feature mask via sparsemax, selecting a handful of features and ignoring the rest; the L-TAE \cite{ltae2020}, which applies a single learned query over the temporal sequence; a standard multi-head Transformer encoder \cite{vaswani2017, russwurm2020}; and TempCNN \cite{tempcnn2019}. Detailed equations for these standard models, and the field-level hybrid-voting rule, are in the Supplementary Material (Sec.~S-A). The Transformer is included as a deliberate \emph{control}: with three pre-norm multi-head self-attention layers it is the most expressive attention model we test---every time step attends to every other---yet, unlike L-TAE-S below, it contains no feature-sparsity mechanism, so it separates attention \emph{capacity} from feature \emph{sparsity}. Patch-based models (a multi-channel 2D CNN and a 3D CNN) tile each field into 100~m patches resampled to $128\times128$ and are likewise trained as five-seed ensembles; details in the Supplementary Material (Sec.~S-A).

\subsubsection{Sparse-Attention L-TAE (L-TAE-S)}
To test whether a tree-like, sparse inductive bias can be grafted onto a temporal-attention encoder, we introduce L-TAE-S, which inserts a TabNet-inspired channel-selection gate before the attention heads (Fig.~\ref{fig:ltae_s}). For each of the $n_{\text{head}}=16$ heads, a per-head linear projection of the embedded, positionally-encoded sequence ($d_{\text{model}}=128$, key dimension $d_k=8$) is passed through \emph{sparsemax} \cite{sparsemax2016} to produce a sparse, instance-dependent mask over embedding channels. Formally, for head $h$ and time step $\mathbf{z}_t\in\mathbb{R}^{E}$,
\[
  \mathbf{m}_{h,t}=\operatorname{sparsemax}\!\big(\mathbf{P}_{h,t}\odot
    \operatorname{BN}(W_h\mathbf{z}_t)\big),\quad
  \tilde{\mathbf{z}}_{h,t}=\mathbf{m}_{h,t}\odot\mathbf{z}_t,
\]
\[
  \mathbf{P}_{h+1,t}=\mathbf{P}_{h,t}\odot\max(\gamma-\mathbf{m}_{h,t},\,0),
  \qquad \mathbf{P}_{1,t}=\mathbf{1};
\]
head $h$ then forms its keys and values from $\tilde{\mathbf{z}}_{h,t}$ before applying L-TAE attention. Because \emph{sparsemax} projects onto the probability simplex, most channels of $\mathbf{m}_{h,t}$ are exactly zero, and the relaxation prior $\mathbf{P}_{h,t}$ (coefficient $\gamma=1.5$) discourages successive heads from re-selecting the same channels; a penalty $\lambda\sum_{h,t}\lVert\mathbf{m}_{h,t}\rVert_{1}$ ($\lambda=10^{-3}$) adds mild additional pressure. This single block---the only architectural difference from the L-TAE---grafts a tree-like sparse feature-selection bias onto the attention backbone while keeping its temporal modeling, mirroring TabNet's sparsemax masks \cite{tabnet_orig}. We evaluate L-TAE-S at both the pixel and field level as a five-seed ensemble under the same weighted-focal objective and field-aggregation protocol as the other temporal networks.

\begin{figure}[!t]
\centering
\includegraphics[width=\columnwidth]{ltae_s.pdf}
\caption{Architecture of the sparse-attention L-TAE-S. Before the multi-head temporal attention, a TabNet-style sparse feature gate (highlighted) applies a per-head \texttt{sparsemax} mask that selects a small subset of embedding channels, with a relaxation prior ($\gamma=1.5$) discouraging heads from re-selecting the same channels. This block is the only architectural difference from the L-TAE; five seeds are ensembled and aggregated to field level.}
\label{fig:ltae_s}
\end{figure}

\subsection{Performance Metrics}
\label{sec:metrics}
Because the classes are severely imbalanced, plain accuracy is misleading. We adopt macro-F1 (F1$_m$, the unweighted mean of per-class F1) as the primary metric, since it weights every crop equally and penalizes failure on the minority cereals where models differ most. We report three further metrics: Cohen's Kappa; weighted F1; and cross-entropy (Xent), the field-level scoring metric of the AI4FoodSecurity challenge \cite{ai4foodsecurity} from which the dataset is drawn (lower is better, computed on hard label predictions). Formal definitions are in the Supplementary Material (Sec.~S-B).

\subsection{Quantifying Uncertainty}
We quantify two sources of variability. For sampling uncertainty in the finite holdout, we apply a nonparametric bootstrap over the 2{,}417 holdout fields: for each model we resample its per-field predictions with replacement $B=10{,}000$ times and recompute macro-F1, reporting the $2.5$th--$97.5$th percentile interval as a 95\% confidence interval (CI); for a between-model gap we take the difference of the two bootstrap distributions and report its CI and a two-sided $p$-value (treating resamples as independent is conservative). For optimization noise, the per-seed standard deviation of ST macro-F1 across the five ensemble members is $0.015$--$0.037$ for the headline temporal models, which is why we report five-seed ensembles throughout. Both sources are small relative to the ranking effects below: every field-bootstrap CI spans only ${\approx}{\pm}0.02$--$0.03$ macro-F1, so we treat ST macro-F1 differences below ${\approx}0.05$ as not statistically distinguishable.

\section{Results}
\label{sec:results}
We evaluate every model under two protocols---in-region field-wise validation and spatial transfer (ST) to the disjoint holdout tile---with all predictions aggregated to the field level.

\subsection{In-Region Validation Overstates Performance and Misranks Models}
Table~\ref{tab:gap} reports both protocols for all models, ranked by ST macro-F1, with the gap $\Delta=\text{F1}_m^{\text{ST}}-\text{F1}_m^{\text{in-region}}$. Two facts stand out. First, \emph{every} model loses accuracy under transfer, but unevenly: from $-0.04$ (field-level L-TAE-S) to $-0.26$ (CNN-BiLSTM at the pixel level). Second, this unevenness reorders the models. In-region the dense temporal pixel nets cluster at the top (L-TAE 0.78, Transformer 0.78, L-TAE-S 0.77, TempCNN 0.77, CNN-BiLSTM 0.72); a practitioner choosing by in-region F1 would pick from this cluster. On the holdout the cluster fractures: the two in-region leaders (L-TAE pixel and the Transformer, $\Delta=-0.20$) drop to $0.58$, TempCNN to $0.56$, and CNN-BiLSTM ($\Delta=-0.26$) to $0.46$, among the lowest absolute scores of any model.

The winners under transfer are a different group. The best single model is the TabNet pixel ensemble (F1$_m=0.60$, $\kappa=0.54$), which also posts the lowest holdout cross-entropy ($4.42$) of any model, tied on macro-F1 by both sparse-attention L-TAE-S variants ($0.60$). The most \emph{robust} models, those with the smallest gaps, are the sparse learners: field-level L-TAE-S posts the single smallest gap ($\Delta=-0.04$), just ahead of logistic regression ($-0.05$) and the pixel-level tree learners. That L-TAE-S---a deep temporal attention network whose only twist is a learned gate restricting each head to a small channel subset---sits atop the robustness ranking is itself evidence for our thesis. The field bootstrap confirms which differences are real: the best transferers beat CNN-BiLSTM by a wide, significant margin (TabNet over CNN-BiLSTM pixel, $\Delta=0.14$, 95\% CI $[0.11,0.18]$, $p<0.001$), so the inversion is not holdout noise; at the very top the leaders (TabNet and both L-TAE-S variants) are statistically indistinguishable ($0.59$--$0.60$; TabNet vs.\ L-TAE-S, $p=0.80$). We therefore do not claim a single best model; we claim that a sparse-feature-selection group occupies the top of the ST ranking while the dense temporal nets prized by in-region validation fall significantly below it.

\begin{table*}[!t]
\centering
\caption{In-region versus spatial-transfer (ST) performance of all twenty-three models, ranked by ST macro-F1. In-region uses field-wise validation; ST uses the disjoint holdout tile (\texttt{34S\_20E\_259N}). $\Delta$ is the gap; the 95\% CI is the percentile bootstrap on ST macro-F1 from $10{,}000$ resamples of the 2{,}417 holdout fields. Xent is the AI4FoodSecurity field-level scoring metric (lower is better). \emph{Features} marks the automated \texttt{xr\_fresh} statistics (174 features) versus raw reflectance. Bold marks the best value per column.}
\label{tab:gap}
\scriptsize
\renewcommand{\arraystretch}{1.1}
\setlength{\tabcolsep}{4pt}
\begin{tabular}{lllccccccc}
\toprule
 & & & \textbf{In-reg} & & \multicolumn{5}{c}{\textbf{Spatial transfer (ST)}} \\
\cmidrule(lr){4-4}\cmidrule(lr){6-10}
\textbf{Model} & \textbf{Level} & \textbf{Features} & \textbf{F1$_m$} & \textbf{$\Delta$} & \textbf{F1$_m$} & \textbf{95\% CI} & \textbf{$\kappa$} & \textbf{wF1} & \textbf{Xent} \\
\midrule
TabNet (raw)          & pixel & raw & 0.73 & $-0.13$ & \textbf{0.60} & .58--.63 & \textbf{0.54} & 0.70 & \textbf{4.42} \\
L-TAE-S               & field & raw & 0.64 & $\mathbf{-0.04}$ & 0.60 & .58--.62 & 0.52 & 0.69 & 4.77 \\
L-TAE-S               & pixel & raw & 0.77 & $-0.17$ & 0.60 & .57--.62 & 0.53 & \textbf{0.70} & 4.63 \\
L-TAE                 & field & raw & 0.67 & $-0.07$ & 0.59 & .57--.62 & 0.47 & 0.67 & 5.51 \\
L-TAE                 & pixel & raw & \textbf{0.78} & $-0.20$ & 0.58 & .55--.60 & 0.49 & 0.68 & 5.20 \\
Transformer           & pixel & raw & \textbf{0.78} & $-0.20$ & 0.58 & .55--.60 & 0.50 & 0.69 & 4.94 \\
XGBoost               & field & \texttt{xr\_fresh} & 0.70 & $-0.13$ & 0.56 & .54--.59 & 0.47 & 0.68 & 4.89 \\
TempCNN               & pixel & raw & 0.77 & $-0.21$ & 0.56 & .54--.59 & 0.49 & 0.68 & 5.09 \\
Base LightGBM         & pixel & \texttt{xr\_fresh} & 0.62 & $-0.06$ & 0.56 & .54--.59 & 0.47 & 0.66 & 4.85 \\
Logistic Regression   & pixel & \texttt{xr\_fresh} & 0.61 & $-0.05$ & 0.56 & .54--.59 & 0.47 & 0.67 & 4.93 \\
LightGBM              & field & \texttt{xr\_fresh} & 0.68 & $-0.12$ & 0.56 & .53--.58 & 0.45 & 0.66 & 4.99 \\
Voting                & pixel & \texttt{xr\_fresh} & 0.65 & $-0.09$ & 0.56 & .53--.58 & 0.48 & 0.66 & 4.85 \\
Random Forest         & pixel & \texttt{xr\_fresh} & 0.65 & $-0.10$ & 0.55 & .52--.57 & 0.45 & 0.65 & 5.07 \\
Base XGBoost          & pixel & \texttt{xr\_fresh} & 0.64 & $-0.09$ & 0.54 & .52--.57 & 0.47 & 0.66 & 5.01 \\
TempCNN               & field & raw & 0.67 & $-0.14$ & 0.54 & .51--.56 & 0.43 & 0.64 & 6.07 \\
Stacking              & field & \texttt{xr\_fresh} & 0.64 & $-0.11$ & 0.53 & .51--.56 & 0.48 & 0.67 & 5.32 \\
CNN-BiLSTM            & field & raw & 0.64 & $-0.14$ & 0.51 & .48--.53 & 0.33 & 0.54 & 7.86 \\
3D CNN                & patch & raw & 0.70 & $-0.20$ & 0.50 & .48--.53 & 0.29 & 0.54 & 7.98 \\
SMOTE Stacked         & field & \texttt{xr\_fresh} & 0.69 & $-0.21$ & 0.49 & .46--.51 & 0.43 & 0.63 & 5.37 \\
CNN-BiLSTM            & pixel & raw & 0.72 & $-0.26$ & 0.46 & .44--.48 & 0.32 & 0.56 & 7.73 \\
Multi-Channel 2D CNN  & patch & raw & 0.63 & $-0.18$ & 0.45 & .43--.47 & 0.30 & 0.54 & 7.62 \\
TabNet (raw, field-avg) & field & raw & 0.62 & $-0.19$ & 0.43 & .41--.46 & 0.31 & 0.56 & 7.35 \\
TabNet (\texttt{xr\_fresh}) & field & \texttt{xr\_fresh} & 0.51 & $-0.15$ & 0.37 & .34--.40 & 0.25 & 0.53 & 6.13 \\
\bottomrule
\end{tabular}
\end{table*}

\subsection{Pixel Versus Field Level: Aggregation as Variance Reduction}
\label{sec:pixfield}
Several architectures were run at both the pixel and the field level (averaging per-pixel sequences within a field before classification). Table~\ref{tab:pixfield} shows that field-level aggregation acts as a variance-reduction regularizer for the \emph{dense} temporal networks: L-TAE's gap shrinks from $-0.20$ to $-0.07$, and CNN-BiLSTM, the worst pixel-level transferer, recovers most, its gap nearly halving from $-0.26$ to $-0.14$ (F1$_m$ $0.46\rightarrow0.51$). The sparse-attention L-TAE-S follows the same dense-backbone pattern despite its channel gating: its gap shrinks from $-0.17$ to $-0.04$ (the smallest of any model) while ST F1$_m$ holds at $0.60$, so the temporal backbone, not the gate, governs the response to aggregation. TabNet is the informative exception: aggregating its raw inputs to the field level \emph{hurts} (F1$_m$ $0.60\rightarrow0.43$), because its robustness already derives from sparse feature selection and field-level training starves it of the per-pixel observation volume it exploits. The tree learners transfer well at either level, and their pixel$\rightarrow$field change is essentially neutral under transfer (XGBoost $0.54\rightarrow0.56$). In short, dense models need aggregation to generalize; sparse, tree-like models do not. Two further analyses---data efficiency under transfer, and the marginal value of the \texttt{xr\_fresh} representation---are summarized in the Discussion and detailed in the Supplementary Material (Sec.~S-E, S-F).

\begin{table}[!t]
\centering
\caption{Pixel- versus field-level spatial transfer for architectures run at both granularities, ranked by field-level ST F1$_m$. Field aggregation helps the dense temporal nets, is roughly neutral for the tree learners, and degrades TabNet. Bold marks the best value per column.}
\label{tab:pixfield}
\renewcommand{\arraystretch}{1.1}
\setlength{\tabcolsep}{3.5pt}
\footnotesize
\begin{tabular}{lccccl}
\toprule
 & \multicolumn{2}{c}{\textbf{Pixel}} & \multicolumn{2}{c}{\textbf{Field}} & \\
\cmidrule(lr){2-3}\cmidrule(lr){4-5}
\textbf{Architecture} & \textbf{F1$_m$} & \textbf{$\Delta$} & \textbf{F1$_m$} & \textbf{$\Delta$} & \textbf{Field in.} \\
\midrule
L-TAE-S    & \textbf{0.60} & $-0.17$         & \textbf{0.60} & $\mathbf{-0.04}$ & raw-avg \\
L-TAE      & 0.58          & $-0.20$         & 0.59          & $-0.07$         & raw-avg \\
XGBoost    & 0.54          & $-0.09$         & 0.56          & $-0.13$         & \texttt{xr\_fresh} \\
LightGBM   & 0.56          & $\mathbf{-0.06}$ & 0.56          & $-0.12$         & \texttt{xr\_fresh} \\
TempCNN    & 0.56          & $-0.21$         & 0.54          & $-0.14$         & raw-avg \\
CNN-BiLSTM & 0.46          & $-0.26$         & 0.51          & $-0.14$         & raw-avg \\
TabNet     & \textbf{0.60} & $-0.13$         & 0.43          & $-0.19$         & raw-avg \\
\bottomrule
\end{tabular}
\end{table}

\subsection{Per-Class Performance and Persistent Confusions}
Figure~\ref{fig:perclass} breaks accuracy down crop by crop on the holdout for four representative models: the best transferer (TabNet), our sparse-attention L-TAE-S, its dense counterpart L-TAE, and CNN-BiLSTM (the largest faller). The errors are highly uneven. All three robust models classify the dominant lucerne/medics and the distinctive canola well (TabNet F1 0.88 and 0.81; L-TAE-S 0.87 and 0.76; L-TAE 0.84 and 0.66), but every model struggles with the cereals and small-grain grazing: even TabNet reaches only 0.55 on barley, 0.48 on wheat, and 0.30 on small-grain grazing. Adding the sparsity gate (L-TAE $\rightarrow$ L-TAE-S) on the same backbone improves the easier classes (canola $0.66\rightarrow0.76$) while leaving the hard cereals essentially unchanged (barley $0.50\rightarrow0.50$, wheat $0.48\rightarrow0.47$). CNN-BiLSTM, by contrast, loses accuracy on the easy classes too (lucerne/medics $0.71$, canola $0.65$), showing its failure is a general inability to transfer, not a minority-class problem alone. The wheat--barley confusion is consistent with their spectral overlap and appears across every model.

\begin{figure}[!t]
\centering
\includegraphics[width=\columnwidth]{per_class_f1.pdf}
\caption{Per-class F1 on the holdout region for four representative models (TabNet, L-TAE-S, L-TAE, and CNN-BiLSTM). TabNet, L-TAE-S, and L-TAE classify lucerne/medics and canola well; the cereals (barley, wheat) and small-grain grazing concentrate the errors. CNN-BiLSTM's accuracy drops on every class. Error bars are 95\% bootstrap CIs over holdout fields.}
\label{fig:perclass}
\end{figure}

\section{Discussion}
\label{sec:discussion}
The conclusion about which model is ``best'' depends almost entirely on whether evaluation is in-region or under spatial transfer, and a model's transfer behavior is predictable from its inductive bias.

\subsection{Why Tree-Like Inductive Bias Transfers}
The clearest pattern is that transfer tracks \emph{inductive bias}, not model family or capacity per se. Figure~\ref{fig:bias} groups every model by inductive bias and plots its ST gap: linear models lose the least (mean $\Delta=-0.05$), tree-based models little (mean $-0.11$), the sparse-attention L-TAE-S sits essentially with the trees (mean $-0.11$ across its variants), while dense temporal/patch nets lose the most (mean $-0.21$) and SMOTE-augmented stacking is similarly fragile ($-0.21$). RF and gradient-boosted trees have been the remote-sensing workhorse for two decades not because they top in-region accuracy (dense nets often do) but because they survive when a model is applied to an unseen region \cite{belgiu2016, maxwell2019}---the operational case. A tree decides one feature at a time, choosing the few features that matter for each split and ignoring the rest; that sparsity is precisely what keeps it from latching onto region-specific quirks of phenology, soil, and atmosphere that pull richer models off course. Independent corroboration spans settings: a 3D-CNN collapses under cross-sensor hyperspectral transfer while a sparse SVM holds \cite{pankajakshan2025}; RF matches or exceeds a dense CNN-LSTM on label-scarce Ghana fields \cite{rustowicz2019}; a maize model drops 16 points moving across East Africa \cite{jin2019}; and bagged/boosted trees out-transfer an elaborate instance-reweighting scheme from the U.S.\ to northwest China \cite{mai2025}. TabNet, the best transferer here, is a neural network built to behave like a tree: a small learned gate picks a handful of features at each step \cite{tabnet_orig}, so it clusters with the trees and inherits their transferability. By contrast, CNN-BiLSTM, TempCNN, L-TAE, and the Transformer consume the whole sequence at once with no feature selection, fit the training region tightly, and transfer poorly---consistent with evidence that the dates such models attend to are themselves contingent on the training region \cite{obadic2022}.

The Transformer is a particularly clean control: although it deploys the most expressive attention of any model we test, far richer than L-TAE's single query, it transfers no better than L-TAE (F1$_m=0.58$, $\Delta=-0.20$). Adding attention \emph{capacity} thus does nothing for transfer, whereas adding feature \emph{selection} does: L-TAE-S and TabNet, the two models that explicitly select features before computing with them, are the best transferers. This isolates feature selection---not the attention mechanism, the temporal backbone, or model size---as the property governing transfer. The L-TAE-S experiment adds that this is \emph{not} a tree-versus-deep tradeoff: the same selection step that drives tree splits can be grafted onto a dense architecture at the cost of a single block, and the payoff is large---L-TAE-S posts the smallest gap in the benchmark and ties TabNet for the best holdout accuracy. The next generation of operational crop classifiers need not choose between the expressivity of deep learning and the transfer-robustness of trees.

\begin{figure}[!t]
\centering
\includegraphics[width=\columnwidth]{inductive_bias_gap.pdf}
\caption{Spatial-transfer gap ($\Delta$ macro-F1) grouped by inductive-bias family. Sparse/tree-like and linear models transfer (small $|\Delta|$); dense temporal/patch nets and synthetic augmentation overfit the training region. L-TAE-S, which adds sparse feature selection to a dense attention backbone, tracks the tree-based models. Bars are family means; points are individual models.}
\label{fig:bias}
\end{figure}

\subsection{Two Robustness Levers}
Two caveats sharpen rather than weaken the story, and together they map the design space. First, it is \emph{pixel-level} TabNet that wins; the field-aggregated TabNet variants are among the worst transferers (Table~\ref{tab:gap}), because aggregation removes the per-pixel observation volume the sparse model exploits. Second, dense temporal attention transfers well \emph{when} applied at the field level: L-TAE's pixel-level gap shrinks substantially under aggregation, and L-TAE-S ends with an even smaller gap, the smallest of any model (Table~\ref{tab:pixfield}). These point to two distinct, partly substitutable pathways. The first lever is built-in feature selection---the model itself picking a few features at a time (trees, TabNet, and now L-TAE-S); this is \emph{intrinsic} robustness paid for at design time. The second is variance reduction imposed from outside via field-level aggregation, which averages away region-specific pixel noise before the classifier sees it---\emph{external} robustness available even when the architecture is fixed. A model with the first lever gains nothing, and may lose observation volume, from the second; this is exactly why TabNet prefers pixels. But the two compose: L-TAE-S adds a sparse gate to a temporal-attention backbone and then benefits from field aggregation, ending with the smallest transfer gap in the benchmark. Practitioners can read this as a flowchart: if free to choose or modify the model, build feature selection in (lever 1); if constrained to an off-the-shelf dense architecture, aggregate predictions to the field level (lever 2). Two further controlled comparisons---ensemble size and the imbalance-handling objective---shift absolute scores but never reorder the robust and fragile groups (Supplementary Material, Sec.~S-D), confirming the ranking is not an artifact of ensembling or imbalance handling.

\subsection{Automated Features, Patch Models, and Persistent Difficulty}
The classical pipelines paired with \texttt{xr\_fresh} were not only robust but cheap, transferring as well as the best deep models (ST F1$_m$ 0.55--0.56) while training in minutes rather than hours (Supplementary Material, Sec.~S-C). A controlled comparison holding the learner fixed shows the \texttt{xr\_fresh} representation yields no measurable transfer-accuracy gain over raw reflectance for the trees; its value is operational (a compact, fast, agronomically interpretable vector) rather than an accuracy edge (Sec.~S-F). This is, to our knowledge, the first evaluation of \texttt{xr\_fresh} in crop classification. Patch-based convolutional models did not exceed pixel-level models and transferred poorly (3D CNN F1$_m=0.50$; multi-channel 2D CNN $0.45$): at 10~m resolution the native 100~m patches are only ${\sim}10\times10$ pixels, so the spatial context is limited and resampling adds interpolation rather than detail, and a dense convolutional representation offers no transfer-protective sparsity. Finally, small-grain grazing, barley, and wheat are the hardest classes for every model (Fig.~\ref{fig:perclass}) and account for most accuracy lost under transfer; that even robust models struggle here, and that CNN-BiLSTM loses accuracy on every class, indicates the limitation is partly intrinsic spectral similarity and partly the transfer problem itself, not class imbalance alone. We also find the most transferable models are the most data-efficient, retaining near-full holdout accuracy on as little as 25\% of the training fields (Supplementary Material, Sec.~S-E).

\subsection{Limitations}
The study uses a single growing season (2017) and a single holdout tile that is spatially \emph{adjacent} to the training tiles within the same agroecological zone. This is a relatively local shift, so the gaps we report most likely \emph{understate} the degradation that fully cross-regional or cross-year transfer would produce; our numbers are best read as a conservative lower bound, and multi-region, multi-year holdouts are needed to confirm the inductive-bias ordering is stable under larger shifts. The labels derive from a public dataset rather than our own campaign, and we use Sentinel-2 optical data only---a deliberate scoping choice that nonetheless limits the achievable ceiling for optically similar cereals, making SAR fusion a clear next step.

\section{Conclusion}
In-region validation, the protocol the crop-classification literature reports almost universally, misranks twenty-three Sentinel-2 classifiers when they are applied to a spatially disjoint holdout. Dense temporal networks that lead under conventional evaluation fall sharply under transfer, the worst of them (CNN-BiLSTM) landing near the bottom, while parsimonious models such as tree ensembles, logistic regression, and a TabNet ensemble retain rank and accuracy. The mechanism is the same property that made tree ensembles the remote-sensing default for two decades: a sparse, axis-aligned feature-selection bias that commits to a few robust signals and ignores the rest. TabNet inherits it through sparsemax masks; a Transformer with far greater attention capacity does not, and transfers no better than the lightweight L-TAE, pointing at sparsity, not capacity, as the lever. Grafting that lever onto an attention model (L-TAE-S) ties TabNet for the best transfer and, at the field level, posts the smallest in-region-to-holdout gap of any model. Field-level aggregation is a second, partly substitutable lever that composes with sparsity in L-TAE-S but hurts TabNet. We therefore recommend: evaluate on a spatially disjoint holdout before declaring one architecture better than another; prefer models with a built-in feature-selection bias when mapping unseen terrain; and, if dense temporal models are required, aggregate their predictions to the field level. The conditions that make the Western Cape difficult are the conditions under which agricultural remote sensing must operate across much of the Global South, and evaluating models the way they are actually deployed---applied to regions where no labels were collected---is what will make remote sensing deliver where it matters most.

\section*{Data Availability}
Ground-truth crop labels were obtained from Radiant Earth's MLHub \cite{mlhub}; Sentinel-2 imagery was extracted via Google Earth Engine. The processed dataset is available from the corresponding author on reasonable request.

\section*{Conflict of Interest}
The authors declare that they have no competing interests.

\begin{thebibliography}{99}
\bibitem{begue2020}
Bégué A, Leroux L, Soumaré M, Faure JF, Diouf AA, Augusseau X, Touré L, Tonneau JP (2020) Remote sensing products and services in support of agricultural public policies in Africa: overview and challenges. \textit{Front Sustain Food Syst} 4:58. \href{https://doi.org/10.3389/fsufs.2020.00058}{https://doi.org/10.3389/fsufs.2020.00058}
\bibitem{ieee2017}
Ienco D, Gaetano R, Dupaquier C, Maurel P (2017) Land cover classification via 
multitemporal spatial data by recurrent neural networks. \textit{arXiv preprint} 
arXiv:1704.04055. \href{https://arxiv.org/abs/1704.04055}{paper}
\bibitem{mahlayeye2024}
Mahlayeye M, Darvishzadeh R, Nelson A (2024) Characterising maize and intercropped maize spectral signatures for cropping pattern classification. \textit{Int J Appl Earth Obs Geoinf} 128:103699. \href{https://doi.org/10.1016/j.jag.2024.103699}{https://doi.org/10.1016/j.jag.2024.103699}
\bibitem{mann2023}
Mann ML, Colson L, Nealon R, Engstrom R, Nakacwa S (2023) Lite Learning: Efficient Crop Classification in Tanzania Using Traditional Machine Learning \& Crowd Sourcing. \textit{SSRN Preprint}.
\bibitem{mann2017}
Mann ML, Warner JM (2017) Ethiopian wheat yield and yield gap estimation: A spatially explicit small area integrated data approach. \textit{Field Crops Res} 201:60--74. \href{https://doi.org/10.1016/j.fcr.2016.10.014}{https://doi.org/10.1016/j.fcr.2016.10.014}
\bibitem{mann2019}
Mann ML, Warner JM, Malik AS (2019) Predicting high-magnitude, low-frequency crop losses using machine learning: an application to cereal crops in Ethiopia. \textit{Clim Change} 154:211--227. \href{https://doi.org/10.1007/s10584-019-02432-7}{https://doi.org/10.1007/s10584-019-02432-7}
\bibitem{maxwell2019}
Maxwell AE, Warner TA, Fang F (2018) Implementation of machine-learning classification in remote sensing: an applied review. \textit{Int J Remote Sens} 39(9):2784--2817. \href{https://doi.org/10.1080/01431161.2018.1433343}{https://doi.org/10.1080/01431161.2018.1433343}
\bibitem{belgiu2016}
Belgiu M, Dr\u{a}gu\c{t} L (2016) Random forest in remote sensing: A review of applications and future directions. \textit{ISPRS J Photogramm Remote Sens} 114:24--31. \href{https://doi.org/10.1016/j.isprsjprs.2016.01.011}{https://doi.org/10.1016/j.isprsjprs.2016.01.011}
\bibitem{zhu2017}
Zhu XX, Tuia D, Mou L, et al. (2017) Deep learning in remote sensing: A comprehensive review and list of resources. \textit{IEEE Geosci Remote Sens Mag} 5(4):8--36. \href{https://doi.org/10.1109/MGRS.2017.2762307}{https://doi.org/10.1109/MGRS.2017.2762307}
\bibitem{russwurm2018}
Ru\ss{}wurm M, K\"orner M (2018) Multi-Temporal Land Cover Classification with Sequential Recurrent Encoders. \textit{ISPRS Int J Geo-Inf} 7(4):129. \href{https://doi.org/10.3390/ijgi7040129}{https://doi.org/10.3390/ijgi7040129}
\bibitem{tempcnn2019}
Pelletier C, Webb GI, Petitjean F (2019) Temporal Convolutional Neural Network for the Classification of Satellite Image Time Series. \textit{Remote Sens} 11(5):523. \href{https://doi.org/10.3390/rs11050523}{https://doi.org/10.3390/rs11050523}
\bibitem{pankajakshan2025}
Pankajakshan P, Padmasanan A, Sundar S (2025) Deep Architectures Fail to Generalize: A Lightweight Alternative for Agricultural Domain Transfer in Hyperspectral Images. \textit{Sensors} 26(1):174. \href{https://doi.org/10.3390/s26010174}{https://doi.org/10.3390/s26010174}
\bibitem{ltae2020}
Garnot VSF, Landrieu L (2020) Lightweight Temporal Self-attention for Classifying Satellite Images Time Series. In: \textit{Advanced Analytics and Learning on Temporal Data (AALTD 2020), ECML PKDD Workshops}, LNCS vol 12588, pp 171--181. Springer. \href{https://doi.org/10.1007/978-3-030-65742-0_12}{https://doi.org/10.1007/978-3-030-65742-0\_12}
\bibitem{jin2019}
Jin Z, Azzari G, You C, Di Tommaso S, Aston S, Burke M, Lobell DB (2019) Smallholder maize area and yield mapping at national scales with Google Earth Engine. \textit{Remote Sens Environ} 228:115--128. \href{https://doi.org/10.1016/j.rse.2019.04.016}{https://doi.org/10.1016/j.rse.2019.04.016}
\bibitem{rustowicz2019}
Rustowicz RM, Cheong R, Wang L, Ermon S, Burke M, Lobell D (2019) Semantic Segmentation of Crop Type in Africa: A Novel Dataset and Analysis of Deep Learning Methods. In: \textit{Proc IEEE/CVF CVPR Workshops}, pp 75--82.
\bibitem{ma2024}
Ma Y, Chen S, Ermon S, Lobell DB (2024) Transfer learning in environmental remote sensing. \textit{Remote Sens Environ} 301:113924. \href{https://doi.org/10.1016/j.rse.2023.113924}{https://doi.org/10.1016/j.rse.2023.113924}
\bibitem{cropformer2023}
Wang H, Chang W, Yao Y, Yao Z, Zhao Y, Li S, Liu Z, Zhang X (2023) Cropformer: A new generalized deep learning classification approach for multi-scenario crop classification. \textit{Front Plant Sci} 14:1130659. \href{https://doi.org/10.3389/fpls.2023.1130659}{https://doi.org/10.3389/fpls.2023.1130659}
\bibitem{mai2025}
Mai J, Feng Q, Fu S, Wang R, Zhang S, Zhang R, Liang T (2025) Enhancing Crop Type Mapping in Data-Scarce Regions Through Transfer Learning: A Case Study of the Hexi Corridor. \textit{Remote Sens} 17(9):1494. \href{https://doi.org/10.3390/rs17091494}{https://doi.org/10.3390/rs17091494}
\bibitem{mlhub}
Radiant Earth Foundation (2024) South Africa Crop Type Competition. Source Cooperative repository.
\bibitem{primary_repo}
Mann ML (2025) South\_Africa\_Crop\_Comp: Source code for multitemporal crop classification. \textit{GitHub repository}
\href{https://github.com/mmann1123/South_Africa_Crop_Comp}{Github}
\bibitem{breiman2001}
Breiman L (2001) Random forests. \textit{Machine Learning} 45(1):5--32. \href{https://doi.org/10.1023/A:1010933404324}{https://doi.org/10.1023/A:1010933404324}
\bibitem{chen2016}
Chen T, Guestrin C (2016) XGBoost: A scalable tree boosting system. In: \textit{Proc 22nd ACM SIGKDD Int Conf Knowl Discov Data Min}, pp 785--794. \href{https://doi.org/10.1145/2939672.2939785}{https://doi.org/10.1145/2939672.2939785}
\bibitem{ke2017}
Ke G, Meng Q, Finley T, et al. (2017) LightGBM: A highly efficient gradient boosting decision tree. \textit{NeurIPS} 30.
\bibitem{wardlow2007towards}
Wardlow BD, Egbert SL, Kastens JH (2007) Analysis of time-series MODIS 250 m vegetation index data for crop classification in the U.S. Central Great Plains. \textit{Remote Sens Environ} 108(3):290--310. \href{https://doi.org/10.1016/j.rse.2006.11.021}{https://doi.org/10.1016/j.rse.2006.11.021}
\bibitem{cnn_multitemp}
Kussul N, Lavreniuk M, Skakun S, Shelestov A (2017) Deep Learning Classification of Land Cover and Crop Types Using Remote Sensing Data. \textit{IEEE Geosci Remote Sens Lett} 14(5):778--782. \href{https://doi.org/10.1109/LGRS.2017.2681128}{https://doi.org/10.1109/LGRS.2017.2681128}
\bibitem{russwurm2020}
Ru\ss{}wurm M, K\"orner M (2020) Self-attention for raw optical Satellite Time Series Classification. \textit{ISPRS J Photogramm Remote Sens} 169:421--435. \href{https://doi.org/10.1016/j.isprsjprs.2020.06.006}{https://doi.org/10.1016/j.isprsjprs.2020.06.006}
\bibitem{fan2023hcrnn}
Fan X, Chen L, Xu X, Yan C, Fan J, Li X (2023) Land Cover Classification of Remote Sensing Images Based on Hierarchical Convolutional Recurrent Neural Network. Forests 14(9):1881. https://doi.org/10.3390/f14091881
\bibitem{tufail2025cnnrnn}
Tufail R, Tassinari P, Torreggiani D (2025) Deep Learning Applications for Crop Mapping Using Multi-Temporal Sentinel-2 Data and Red-Edge Vegetation Indices: Integrating Convolutional and Recurrent Neural Networks. Remote Sens 17(18):3207. https://doi.org/10.3390/rs17183207
\bibitem{bandar2024bilstm}
Bandar A, Co\c{s}kun\c{c}ay A (2024) Crop Classification with Attention Based BI-LSTM and Temporal Convolution Neural Network Combination for Remote Sensing BreizhCrop Time Series Data. Y\"uz\"unc\"u Y{\i}l \"Universitesi Fen Bilimleri Enstit\"us\"u Dergisi 29(1):173--188. https://doi.org/10.53433/yyufbed.1335866
\bibitem{espejo2020}
Espejo-Garcia B, Mylonas N, Athanasakos L, Fountas S, Vasilakoglou I (2020) Towards weeds identification assistance through transfer learning. \textit{Comput Electron Agric} 171:105306. \href{https://doi.org/10.1016/j.compag.2020.105306}{https://doi.org/10.1016/j.compag.2020.105306}
\bibitem{koklu2022}
Koklu M, Unlersen MF, Ozkan IA, Aslan MF, Sabanci K (2022) A CNN-SVM study based on selected deep features for grapevine leaves classification. \textit{Measurement} 188:110425. \href{https://doi.org/10.1016/j.measurement.2021.110425}{https://doi.org/10.1016/j.measurement.2021.110425}
\bibitem{yang2020}
Yang S, Gu L, Li X, Jiang T, Ren R (2020) Crop Classification Method Based on Optimal Feature Selection and Hybrid CNN-RF Networks for Multi-Temporal Remote Sensing Imagery. \textit{Remote Sens} 12(19):3119. \href{https://doi.org/10.3390/rs12193119}{https://doi.org/10.3390/rs12193119}
\bibitem{mann2025}
Mann ML (2025) xr\_fresh: Automated Time Series Feature Extraction for Remote Sensing and Gridded Data. \textit{J Open Source Softw} 10(115):9009. \href{https://doi.org/10.21105/joss.09009}{https://doi.org/10.21105/joss.09009}
\bibitem{ai4foodsecurity}
AI4EO (2021) AI4FoodSecurity Challenge (South Africa), European Space Agency (ESA), in partnership with Planet, TUM, DLR and Radiant Earth. \href{https://platform.ai4eo.eu/ai4food-security-south-africa}{https://platform.ai4eo.eu/ai4food-security-south-africa}
\bibitem{james2013}
James G, Witten D, Hastie T, Tibshirani R (2013) \textit{An Introduction to Statistical Learning}. Springer.
\bibitem{friedman2001}
Friedman JH (2001) Greedy function approximation: A gradient boosting machine. \textit{Ann Stat} 29(5):1189--1232. \href{https://doi.org/10.1214/aos/1013203451}{https://doi.org/10.1214/aos/1013203451}
\bibitem{penabarragan2011}
Pe\~na-Barrag\'an JM, Ngugi MK, Plant RE, Six J (2011) Object-based crop identification using multiple vegetation indices, textural features and crop phenology. \textit{Remote Sens Environ} 115:1301--1316. \href{https://doi.org/10.1016/J.RSE.2011.01.009}{https://doi.org/10.1016/J.RSE.2011.01.009}
\bibitem{tabnet_orig}
Arik SÖ, Pfister T (2021) TabNet: Attentive Interpretable Tabular Learning. \textit{Proc AAAI Conf Artif Intell} 35(8):6679--6687. \href{https://doi.org/10.1609/aaai.v35i8.16826}{https://doi.org/10.1609/aaai.v35i8.16826}
\bibitem{tabnet2024}
Muruganantham P, Wibowo S, Grandhi S, Islam N (2024) A Deep Learning Approach for Dealing with Tabular Data in Crop Classification. \textit{Proc IEEE Int Conf ICT Convergence (ICTC)}:2054--2059. doi:10.1109/ICTC62082.2024.10826760
\bibitem{vaswani2017}
Vaswani A, Shazeer N, Parmar N, Uszkoreit J, Jones L, Gomez AN, Kaiser \L{}, Polosukhin I (2017) Attention Is All You Need. In: \textit{Advances in Neural Information Processing Systems (NeurIPS)}, vol 30, pp 5998--6008. \href{https://arxiv.org/abs/1706.03762}{arXiv:1706.03762}
\bibitem{sparsemax2016}
Martins A, Astudillo R (2016) From Softmax to Sparsemax: A Sparse Model of Attention and Multi-Label Classification. In: \textit{Proceedings of the 33rd International Conference on Machine Learning (ICML)}, PMLR vol 48, pp 1614--1623. \href{https://proceedings.mlr.press/v48/martins16.html}{paper}
\bibitem{obadic2022}
Obadic I, Roscher R, Oliveira DAB, Zhu XX (2022) Exploring Self-Attention for Crop-type Classification Explainability. \textit{arXiv preprint} arXiv:2210.13167. \href{https://arxiv.org/abs/2210.13167}{https://arxiv.org/abs/2210.13167}
\end{thebibliography}


\end{document}
